To measure the contribution of IV dynamics, we compared five variants on identical stock paths for the four scenario contracts. The comparison held contracts, entry premiums, and pricing assumptions fixed so that differences in interim marks could be attributed to the IV module (Section~\ref{sec:ablation_methods}). After ten trading transitions, the full factor changed marks by a mean absolute \$1.77--\$2.41 per share relative to direct surface evaluation (Supplementary Table~\ref{tab:dynamic_ablation}). These pathwise changes were much larger than the corresponding difference in the GS put's mean P\&L, which was only \$0.05 lower under the full factor with a paired Monte Carlo standard error of \$0.04. Positive and negative differences therefore largely canceled when averaged across paths. Frozen IV and deterministic mean reversion also produced different marks from the direct surface, showing that movement of the target and persistence around it affected valuation even without stochastic shocks. Comparing the coupled and uncoupled factors on the same paths showed how return coupling changed valuation across stock outcomes. For both tickers, short-put P\&L declined most on falling-stock paths, while short-call P\&L improved most on rising-stock paths (Fig.~\ref{fig:dynamic_ablation}). Coupling also changed the lower tails of short-position P\&L. We measured those tails with 5\% expected shortfall, defined as mean P\&L at or below its 5\% quantile, for which more negative values indicate larger simulated losses. Negative coupling moved the GS put's expected shortfall from $-\$63.24$ under uncoupled factors to $-\$64.86$ under the full factor. The call's expected shortfall improved from $-\$71.54$ to $-\$69.83$, and the LLY contracts showed the same direction of change (Supplementary Table~\ref{tab:dynamic_ablation}). The full-factor mean absolute mark differences were similar across the three evaluation seeds (Supplementary Table~\ref{tab:ablation_seeds}). All variants produced identical entry marks and terminal payoffs. The experiment thus located the contribution of IV dynamics in interim valuation and liquidation risk while preserving the stock paths and terminal outcomes.

\begin{figure}[H]
\centering
\includegraphics[width=\textwidth]{sections/figures/dynamic_ablation.pdf}
\caption{Effect of return coupling on interim short-option P\&L as a function of stock return after ten trading transitions for (a) GS and (b) LLY. Each difference is coupled-factor P\&L ($\rho=-0.6$) minus uncoupled-factor P\&L ($\rho=0$) on the same stock path, with the same independent normal innovations. Negative values indicate worse short-position P\&L under coupling. The 3{,}000 paths per ticker are sorted by stock return and divided into ten groups of 300. Symbols show median paired P\&L differences at median stock returns; connecting lines guide the eye, and shaded bands span the within-group 25th--75th percentiles. These bands describe variation across paths, not confidence intervals. Red circles denote puts and navy diamonds denote calls. All marks are evaluated on May 12, 2026, with 17 calendar days to expiry.}\label{fig:dynamic_ablation}
\end{figure}

We examined numerical accuracy and strike consistency to determine how the ablation differences should be interpreted as option prices. Increasing lattice depth from 201 to 401 changed sampled marks by at most \$0.0017 per share, well below the differences among IV variants. No factor reached the variance floor in the four scenario contracts. Current strike/spot nevertheless left the training interval on 0.7--2.9\% of pre-expiry path dates across ticker, contract, and seed combinations, so some marks used extrapolated surface values. The strike audit then tested necessary conditions across separately priced contracts. At 401 lattice steps, the direct GS surface exceeded the \$0.01 convexity tolerance in 178 of 1{,}818 tested strike triples (Supplementary Table~\ref{tab:strike_audit}). Its largest excess was \$0.113 per share. The full factor had three monotonicity violations among 2{,}020 adjacent-strike comparisons, with a maximum of \$0.013. These violations persisted at the higher lattice depth, while no LLY variant violated the tested conditions at the stated tolerance. The checks therefore distinguished two issues in the calculation. The sampled ablation effects exceeded the changes from numerical refinement, but single-contract lattice accuracy did not ensure consistency across strikes. The residual GS violations limited the interpretation of the generated prices as market-consistent training data even though the same outputs supported controlled comparisons of IV assumptions.
\FloatBarrier
